\chapter{Bestehende Ansätze}
\label{ch:ansatz}


\section{EU Regulierung und Framework zur Digitalsouveränität und Cloud}

\subsection{DSGVO} \cite{gdpr2016}
Um die Datensouveränität ihre Bevölkerung innerhalb der EU zu schaffen hat in Mai 2018 die EU die Datenschutz-Grundverordnung (DSGVO) eingeführt.
Die DSGVO ist eine Verordnung der EU, die die Verarbeitung personenbezogener Daten regelt und darauf abzielt, die Privatesphäre von Einzelpersonen zu stärken.
Sie legt Richtlinien für die Erhebung, Verarbeitung, Speicherung und Nutzung personenbezogener Informationen fest und stellt sicher, 
dass Einzelpersonen mehr Kontrolle über ihre Daten haben. 
Zu den wichtigsten Grundsätzen gehören Transparenz, Datenminimierung sowie das Recht auf Zugang zu personenbezogenen Daten und deren Löschung.

\subsection{Gaia-X} \cite{GaiaXOverview2025}
Gaia-X ist ein europäisches Projekt zur Schaffung eines vertrauenswürdigen und interoperablen Daten- und Infrastruktur-Ökosystems. Ziel ist es, Datenhoheit, Sicherheit, Transparenz, Nachhaltigkeit und fairen Wettbewerb in Europa zu fördern und die Abhängigkeit von nicht-europäischen Cloud-Anbietern zu reduzieren. Kernbestandteile sind das Gaia-X Trust Framework, das Regeln, Spezifikationen und Prüfmechanismen definiert, sowie das Gaia-X Digital Clearing House (GXDCH), ein technisches System, das überprüft, ob Organisationen und Dienste die Anforderungen des Trust Frameworks erfüllen. Nach erfolgreicher Prüfung erhalten Organisationen ein Gaia-X Label bzw. verifizierbare digitale Nachweise (Verifiable Credentials), die Transparenz, Datenschutz, Sicherheit, Interoperabilität, Portabilität, Nachhaltigkeit und europäische Kontrolle dokumentieren.

\subsection{Cloud Sovereignty Framework} \cite{EC_CloudSovereigntyFramework_2025}
Cloud Sovereignty Framework ist ein Framework von europäische Kommission, das die für die Bereitstellung von Cloud-Diensten relevanten Souveränitätszielen definiert.
Mit den Objektiven kann man den Cloud Service Provider bewerten, wie gut er die erforderlichen Souveränitätskriterien durch Sovereignty Effectiveness Assurance Level (SEAL) erfüllt. die SEAL Level dient als Minimum Assurance Level, die durch Kundenspezifikation definiert werden. Die Kunden rechnen die Sovereignty Score von dem Cloud Service Provider, um auszusortieren, welche Cloud Service Provider zu ihren Souveränitätskriterien am besten passt.